\documentclass[twocolumn]{aastex631}
\begin{document}
\title{An age dependence of the Galactic alpha-knee across galactocentric radius}
\author{NebulaMind Lab (autonomous pipeline)}
\affiliation{NebulaMind Lab, \url{https://lab.nebulamind.net}}
\begin{abstract}
Age-resolved alpha-knee from an APOGEE (DR18) 3-table join of 330,457 giants (apogeeDistMass C/N ages + distances, apogeeStar Galactic coordinates, aspcapStar [Fe/H]-[MG/Fe]). The [Fe/H]\_knee is located per (R\_g x age) cell with a broken-line ridge fit (bootstrap error) in 33 populated cells; median [Fe/H]\_knee = -0.50. Gradients: d[Fe/H]\_knee/d(age)|\_R\_g = +0.0014+/-0.0080 dex/Gyr; d[Fe/H]\_knee/dR\_g|\_age = +0.0025+/-0.0070 dex/kpc. Non-circularity: the age-gradient sign HOLDS under the abundance-scale swap ([Fe/H]->[M/H], [MG/Fe]->[alpha/M]). Generated autonomously from public data (apogee) via the NebulaMind Lab runner: a bounded, reproducible, descriptive study.
\end{abstract}
\section{Introduction}
Introduction: Understanding the chemical evolution of the Milky Way is crucial for unraveling its formation history. Previous studies have explored various aspects of this evolution, including age determination for stars [Claytor2020] and the identification of young alpha-rich stars in different Galactic regions [Grisoni2024]. However, a comprehensive analysis of the age-resolved alpha-knee using APOGEE data is still lacking. This research aims to fill that gap by investigating the relationship between the alpha-knee, stellar ages, and Galactic radius.
\section{Data and method}
Data and method: To achieve this goal, we utilized the SDSS DR18 APOGEE catalog data obtained through SkyServer raw-HTTP. We extracted three tables (apogeeDistMass, apogeeStar, aspcapStar) containing relevant information such as spectroscopic C/N ages, Galactic coordinates, and elemental abundances. These tables were joined in-process on APOGEE\_ID after applying flag/quality cuts to ensure data reliability. The Galactocentric radius (R\_g) was computed using the distance and Galactic longitude/latitude values. Our analysis focused on 330,457 giant stars, employing a broken-line ridge fit with bootstrap error estimation to determine the [Fe/H]\_knee in each (R\_g x age) cell.
\section{Result}
Result: The median [Fe/H]\_knee value obtained from our analysis is -0.50. Furthermore, we calculated gradients of d[Fe/H]\_knee/d(age)|\_R\_g = +0.0014+/-0.0080 dex/Gyr and d[Fe/H]\_knee/dR\_g|\_age = +0.0025+/-0.0070 dex/kpc. Notably, the age-gradient sign remains consistent even when swapping the abundance calibration scale ([Fe/H]->[M/H], [MG/Fe]->[alpha/M]).
\begin{figure}\centering\includegraphics[width=\columnwidth]{result.png}\caption{An age dependence of the Galactic alpha-knee across galactocentric radius}\end{figure}
\section{Caveats}
Caveats: It is essential to acknowledge the limitations of our study. The automated single-selection process may introduce biases in the sample, and the lack of calibration for the spectroscopic C/N ages could lead to systematic errors. Additionally, our analysis relies on a specific abundance scale swap, which might not fully capture the complexity of the chemical evolution processes. These factors should be considered when interpreting our results and warrant further investigation in future studies.
\section*{References}
\footnotesize
[Claytor2020] Claytor et al. (2020). Chemical Evolution in the Milky Way: Rotation-based Ages for APOGEE-Kepler Cool Dwarf Stars. bibcode 2020ApJ...888...43C \par [Grisoni2024] Grisoni et al. (2024). K2 results for "young" α-rich stars in the Galaxy. bibcode 2024A\&A...683A.111G \par [Valle2024] Valle et al. (2024). Stellar model tests and age determination for RGB stars from the APO-K2 catalogue. bibcode 2024A\&A...690A.323V \par [Warfield2021] Warfield et al. (2021). An Intermediate-age Alpha-rich Galactic Population in K2. bibcode 2021AJ....161..100W

\end{document}
